%%%%%%%%%%%%%%%%%%%%% ch08_cn.tex, Springer monograph; LLM decision layer %%%
\chapter{巡检分析中的数字孪生与大语言模型}
\label{ch:llm-dt}


\section{面向建筑巡检的数字孪生}
\label{sec:ch8-dt}

城市中大量既有建筑在安全与成本压力下持续服役，人工立面巡检仍劳动密集且风险高。因此无人飞行器（UAV）、传感与人工智能进入结构健康监测（SHM）与建筑运维（O\&M）。本书第2--7章已按该管线依次展开：任务与运动规划、三维重建、数据集与缺陷模型，以及把类型化发现挂到可维护资产上的地理信息。前几章交付的是已被测量、配准并空间组织好的建筑；它们尚未完成的，是日常语言下的决策支持。网格上出现检测框，很少等于日常运维结束。经理仍会问：哪些板块共享渗水路径，自上一航次以来发生了什么变化，下周吊篮应覆盖哪些资产，如何起草工程师可审计的维修简报。本章因此与上游技术栈连续，职责上又有区别：它消费前几章产物，并展开其上的语言与决策层。

数字孪生是回答此类问题的常用名称，但其含义宽于逼真三维浏览器。建设与土木基础设施中的孪生研究强调全生命周期信息流、多方协同，以及``可视化''与``可行动语义''之间的缺口~\cite{boje2020,opoku2021,lu2020}。遗产与建筑记录文献同样强调：在几何之外，还需要与诊断挂钩的语义以及可监测的出处~\cite{bruno2018,tsilimantou2020}。面向设施的讨论则补上一点：运维价值出现在孪生接到工作流程之时，而不是只装饰会议幻灯片之时~\cite{chen2023facility,lu2020}。这些工作共同说明：孪生首先是一种组织与交付信息的方式，而不是某一种固定软件品牌。

建筑巡检领域的孪生实践则更接近一条可复述的管线。无人机航次产生影像与位姿；摄影测量与点云处理形成几何；缺陷检测与语义分割给出类型化发现；BIM 或 GIS 配准把这些发现挂到可维护的建筑对象上~\cite{zhang2022,zhang2025,yang2025detreconreg,zhao2024facade,zhao2026physical}。已发表系统在这一链条上的侧重各不相同：有的强调网格与实景模型的可视化浏览，有的强调缺陷数据库与报告生成，有的进一步把立面板块、楼层或设施台账条目作为长期更新的宿主~\cite{zhang2022,chen2023facility,han2025}。它们很少把``聊天''当作核心产物；真正决定能否支撑日常运维的，是几何、检测、身份与出处能否在多次航次之间保持一致，而不是界面是否足够流畅。

大语言模型（LLM）进入这一图景时，角色是语言接口，而不是第二套检测器，也不是现场数据库。现成模型可以解析混合的空间--症状问句、起草面向不同读者的简报，并遵守答案模板。但若单独使用，它们对点名建筑对象与带版本手册仍不可靠：流畅文本可能虚构板块编码、抹平 RGB/红外分歧，或回忆错手册版本~\cite{ji2023,lewis2020}。检索增强生成、结构化输出与工具调用，正是工程上对这类失效的一般回应：把活事实留在权重之外，查询时取回，并对照检索证据校验主张~\cite{lewis2020,shuster2021,yao2023react}。视觉—语言模型（VLM）若使用，角色同样有界：可以帮助叙述已经挂接的影像裁剪，却不应发明宿主编码或改写检测器标签。邻近的 AEC 语言系统已经展示了这些模式的不同片段：BIM 助手改善模型导航~\cite{zheng2023bimassistant}，规范与安全文档 RAG 改善条款查找~\cite{joffe2025buildingcodes,lee2024ragsafety,krutli2025fmrags}，设施管理问答把维护日志与工单语境接入生成~\cite{krutli2025fmrags,buga2025llmfmsql}，IFC 工具智能体与运维机器人则把语言请求接到结构化动作与孪生更新之上~\cite{gao2026ifcagent,luo2025llmdtrobot,hellin2025bimmas}。它们说明 LLM 在土木里已经``能用''，也说明巡检场景还需要更严格的证据纪律。

由此可以得到一条可行的巡检决策支持思路。把航次事实放在带版本的孪生库中；让语言系统检索以建筑对象为中心的记录、引用可打开证据、在证据包偏薄时拒答，并只在人工或 geofence 门禁之后提议行动。在这种安排里，孪生是长期记忆，LLM 是规则下的阅读者与起草者。缺陷的类型化发现仍由感知负责；接地解释与可排程草稿仍由语言负责。历史巡检手艺其实早已指向同一方向：纸质清单强制点名位置，照片日志强制出处，维修工单强制行动字段。数字孪生与 LLM 增加的是尺度与可查询性，并不允许抛弃这些纪律~\cite{zhang2022,chen2023facility,ji2023}。

实践中读者会遇到多种孪生成熟度。\emph{几何孪生}保存网格、点云或 BIM 外壳；\emph{语义孪生}挂接类型化实体与稳定宿主编码；\emph{决策可用孪生}还进一步暴露观测卡片、用于传播的拓扑，以及检索与工具可调用的接口。把聊天模型硬接到第一级，通常只会得到流畅的虚构。本章的可行目标是第三级，并以第二级为最低前提。组织现实也强化这一目标：无人机承包方、顾问、设施团队与资产经理很少同属一人，因此对象身份、出处与邻接关系必须经得起文件交接，而不能每次都打电话问建网格的人~\cite{boje2020,opoku2021,chen2023facility}。

概念路径清楚之后，再看一个具体场景。仓立面巡检结束后，经理询问西立面中层空鼓是否与上方潮湿渍迹相关。网格浏览器可以同时显示两块区域，检测器可以给两块贴标签；只有语言可用的孪生才能装配相邻宿主记录、引用可打开影像的观测条目，并在证据不足时拒答或起草维修简报。这一场景是后文的贯穿例子：孪生存事实，LLM 在引用规则下提议语言，而不是第二几何源。后文示例使用 \texttt{W-3-17}、\texttt{W-4-17} 一类占位编码仅为便于阅读轨迹；实际项目应替换为设施台账或 BIM 宿主中的真实编码。

只有在上述背景清楚之后，本书才需要给出一幅过程图景：数字孪生与大语言模型如何落在同一条巡检工作流里。第2--7章已把飞行变成几何、类型化缺陷与空间挂接对象；图~\ref{fig:ch8-framework} 说明这些产物如何成为语言可用的孪生记忆、LLM 如何在检索与引用规则下阅读该记忆，以及设施团队如何得到可审计答案、报告与有门禁的维护行动。它是巡检决策支持的过程框架，而不是某一篇论文的软件模块清单。图~\ref{fig:ch8-dt-pipeline} 再给出语言层所消费的建模路径：物理采集、孪生建模与 GIS+BIM 集成如何把缺陷发现变成可检索、可打开证据的记录；此处作为语言分析的输入接口，而不是第二次重建教程~\cite{zhang2022,zhang2025,yang2025detreconreg}。Asset-ID、资产护照与引用或拒答等本章专用术语，在第~\ref{sec:ch8-dt-to-language}节接到第~7章设施身份之后一并定义；在此之前，正文尽量使用``可维护宿主''``设施台账编码''``观测记录''等更中性的说法。

\begin{figure}[htbp]
\centering
\includegraphics[width=\linewidth]{framework_overview_overall.pdf}
\captionsetup{font=normalsize,labelfont=bf,textfont=normalsize}
\caption{无人机建筑巡检中数字孪生与大语言模型的过程框架：航次数据形成可维护孪生；孪生为语言推理提供证据；语言系统向人工运维返回可审计答案、报告与有门禁的维护行动。宜自左向右读作一条巡检工作流，而不是单篇论文的软件栈。}
\captionsetup{font=small,labelfont=bf,textfont=small}
\label{fig:ch8-framework}
\end{figure}

\begin{figure}[htbp]
\centering
\includegraphics[width=\linewidth]{dt_modeling_pipeline.pdf}
\caption{语言层所消费的无人机→孪生建模路径：物理数据采集，涵盖影像、位姿与 BIM 语境；数字孪生建模，涵盖重建、缺陷检测与 BIM 配准，以及把可打开缺陷记录暴露给 GIS+BIM 虚拟集成。此处作为语言分析的输入接口，而不是第二次重建教程。}
\label{fig:ch8-dt-pipeline}
\end{figure}

后文先在第~\ref{sec:ch8-llm-foundations}节说明大语言模型在土木工程里已经如何被使用，再在第~\ref{sec:ch8-dt-to-language}节把前几章产物组织成语言可用的孪生记录与装配规则。本章不引入新的检测器；它说明语言层如何消费第2--7章已经建立的孪生基础。


\section{面向土木工程应用的大语言模型基础}
\label{sec:ch8-llm-foundations}

当代 LLM 建立在 Transformer 架构之上~\cite{vaswani2017attention}。自注意力能够建模长程文本依赖，这对``症状 + 位置 + 规范条款''交织的巡检提问尤为重要。现成模型已能遵循自然语言指令并起草流畅的建筑文本~\cite{brown2020,ouyang2022instructgpt}。思维链提示可引出中间步骤，有助于多跳问题~\cite{wei2022cot}，但并不能保证某个立面板块真的存在于孪生中。施工领域对 GPT 类工具的综述给出同样审慎：能力真实存在，治理与事实接地并不自动成立~\cite{saka2024gptconstruction}。

若只停留在架构与提示，读者仍难判断 LLM 在土木里``究竟怎么用''。近年已发表工作给出了较清晰的应用谱系。\emph{模型导航与查询}方面，BIM 助手与自然语言模型查询帮助用户在复杂模型中定位构件、生成解释或发起 SQL/IFC 查询~\cite{zheng2023bimassistant,guo2025bimllmquery,liu2025bimcoder,li2025b}。\emph{文档与规范检索}方面，RAG 已用于建筑规范、安全手册、设施管理文档与混合法规语料，重点在版本忠实与来源可追溯~\cite{joffe2025buildingcodes,lee2024ragsafety,krutli2025fmrags,jeon2025hybridrag,he2025regulationhybrid}。\emph{结构化动作与孪生交互}方面，IFC 工具智能体、图孪生 GPT 与运维机器人把语言请求接到类型化工具、批准门与孪生更新之上~\cite{gao2026ifcagent,pan2026graphdtgpt,luo2025llmdtrobot,hellin2025bimmas,quan2026bimfacrag}。\emph{设施问答与工单语境}方面，FM 导向 RAG 与 SQL 代理尝试把维护日志、资产台账与工单字段接入回答~\cite{krutli2025fmrags,buga2025llmfmsql,chen2023facility}。这些系统并不等价，却共同说明：土木里的 LLM 价值主要来自``接对数据、接对工具、接对流程''，而不是单独换一个更大的基座模型~\cite{saka2024gptconstruction,ma2026tenquestions}。

除架构与提示之外，土木部署还要分清三种 LLM 行为。遵从指令的模型可以在编排器给出模板时按字段作答，从而在短回复里区分观测与假说~\cite{ouyang2022instructgpt,brown2020}。会调用工具的模型可以请求类型化函数，而不把世界状态写进权重，从而使孪生读取与邻域扩展落在模型之外~\cite{schick2023toolformer,yao2023react}。长上下文模型可以装下打包后的证据摘要，但仍需要校验器：上下文变长并不会创造出从未检索到的对象编码~\cite{lewis2020,ji2023,liang2023helm}。面向施工的 GPT 综述与建筑应用议程因此把治理、数据就绪与人机监督当作与模型选型并列的设计变量~\cite{saka2024gptconstruction,ma2026tenquestions}。对立面巡检，这些变量落到实处，就是版本钉扎、模态标签、引用或拒答检查，以及写入门禁。当代基础模型在接口形态上也不同：对话微调 API 强调遵从指令；工具调用 API 暴露类型化函数 schema；检索增强部署则把现场事实留在权重之外。巡检系统应按契约选择与组合这些接口~\cite{ouyang2022instructgpt,schick2023toolformer,lewis2020,saka2024gptconstruction}。

本章因此采取明确立场：把 LLM 当作\emph{可部署的语言组件}，用提示约束、接到孪生库与工具即可工作；不展开现场专用训练或微调议程。航次事实随飞行更新，对仍在服役的建筑，诚实做法是把这些事实放在权重之外，放入带版本手册与结构化孪生记录，并让模型在查询时读取~\cite{lewis2020,krutli2025fmrags}。还宜分清三种常被混为一谈的行为。\emph{语言能力}是解析与生成流畅建筑文本的能力；\emph{参数回忆}是权重里已有的通用知识，对词汇有用，对具体项目的台账编码不可靠；\emph{接地推理}是依据当前会话检索到或被给予的证据作答的能力。立面运维成败取决于第三种。能写出优美渗水叙事却不引用孪生观测的 LLM，不是巡检助手，只是起草玩具。

本节先介绍面向土木的生成式语言技术，再接到巡检分析。正文不罗列训练配方，而把即用技术组织为可组合的三层决策栈，见表~\ref{tab:ch8-tech-families}。感知仍是第6章的职责；语言技术消费带类型、带模态标签的观测。第 L1 层是\emph{语言接口}：如何用提示与答案模板驾驭现成模型。第 L2 层是\emph{证据接地}：如何用检索、结构化输出与工具把答案钉在孪生记录上。第 L3 层是\emph{记忆与有界流程}：如何用长期孪生记忆与有门智能体支撑多步运维，而不放任无监督行动。下文以连贯叙述展开三层，再说明已发表 AEC 系统如何专门化，或未能专门化，到立面巡检。

\begin{table}[htbp]
\centering
\caption{三层即用 LLM 技术$\rightarrow$立面巡检。}
\label{tab:ch8-tech-families}
\begin{tabularx}{\textwidth}{@{}>{\raggedright\arraybackslash}p{2.2cm}>{\raggedright\arraybackslash}p{3.0cm}X@{}}
\toprule
层 & 核心技术 & 巡检应用 \\
\midrule
L1 语言接口 & 提示、模板、分步解释 & 解析问题；起草；区分观测与假说 \\
L2 证据接地 & RAG、schema、工具/API & 取护照/手册；约束报告字段；引用或拒答 \\
L3 记忆与流程 & 孪生/RAG 记忆、有界智能体 & 权威现场记忆；工单/复访须门禁 \\
\bottomrule
\end{tabularx}
\end{table}

第 L1 层从现成模型在编排器给出细致提示时已经能做的事讲起。模型可以解析``中层空鼓且上层有渍迹''这类空间，症状混合问句，为不同利益相关方起草可读简报，并服从短答案模板：强制引用检索包中已出现的建筑对象，区分观测与假说，检查失败时输出拒答码~\cite{ouyang2022instructgpt,ji2023}。分步提示可以口头化传播假说，但必须在邻域\emph{已经进入}证据包之后；否则模型可能编造听起来很``结构''的故事~\cite{wei2022cot,ji2023}。对立面运维，实用系统提示可以很短：分类意图；只谈检索包中已有的对象与观测；保留第6章模态标签；证据不足时显式拒答并给出原因。少样本示例应演示拒答，而不只演示成功。L1 \emph{不应}被要求用参数回忆记住上一航次台账编码、发明缝长，或在经理催促``直接标成已修''时覆盖孪生事实，那些属于接地与门禁，不属于文笔。

第 L2 层是巡检系统离开聊天演示的地方。检索增强生成在查询时加载带版本手册或观测卡片~\cite{lewis2020,shuster2021}；土木落地已覆盖规范、安全手册、设施管理文档与缺陷对话~\cite{joffe2025buildingcodes,lee2024ragsafety,krutli2025fmrags,jeon2025hybridrag,he2025regulationhybrid}。自适应与自反思变体在检索偏弱时重试~\cite{gao2023hyde,asai2024selfrag,jeong2024adaptive}，但巡检设计选择比``加个 RAG''更尖锐：对资产中心问题，优先实体/护照单元而非任意 PDF 切块~\cite{peng2024}。结构化与约束生成使输出可机检，报告字段、CMMS 映射，以及要求所引 Asset-ID 必须属于证据集的校验~\cite{pan2024b,chen2024bimllmontology,iversen2026,buga2025llmfmsql}。工具与函数调用把推理与世界动作分开：模型可以请求宿主护照、扩展邻域、列出观测、起草报告或提议复访，但必须从孪生索取这些事实，而不是发明它们~\cite{schick2023toolformer,yao2023react,guo2025bimllmquery,liu2025bimcoder,li2025b,li2025c}；孪生写入与调度仍须服务端批准~\cite{luo2025llmdtrobot,gao2026ifcagent}。长上下文不是倾倒整面立面的许可证。打包应优先结构化摘要，Asset-ID、主要症状、近期观测 ID、邻域 ID、短位置短语，并放影像路径而非原始像素，以便人类打开证据；若仍溢出，先丢低置信观测。传播所需的拓扑链接应优先保留~\cite{lewis2020,asai2024selfrag}。

第 L3 层让技术栈跨会话可持续，并在可控条件下可行动。在智能体之前，读者需要记忆阶梯，见表~\ref{tab:ch8-memory}。认知侧常把智能体记忆分为工作、情景、语义与程序性存储；巡检阶梯是这一思想的现场专门化，而不是另起一套分类学。参数记忆是权重里已有的通用知识，对语言有用，对``上周某仓立面裂缝增长''几乎无用~\cite{brown2020}。上下文记忆是提示窗口与短对话状态，即工作记忆：能力强，但受长度与会话限制；若明天新开对话，昨天关于``西侧指装卸台立面''的澄清可能消失，除非外部库保存了别名。检索记忆按需加载带版本手册与观测卡片，即须钉扎版本的情景与语义摘要~\cite{lewis2020,shuster2021,krutli2025fmrags}。结构化长期孪生记忆则是 Asset-ID 库本身，BIM/GIS 中的几何、拓扑与时间戳观测~\cite{pan2024b,zhang2025}。对立面运维，孪生才是权威；模型应通过检索与工具读取它。模型不得假装权重记得现场。有界智能体再把规划、记忆与工具组成工单或复访等流程，并置于 geofence 与人工门之下：程序性记忆住在 harness 里，而不是静默改权重~\cite{wang2024surveyagents,li2024surveymultiagent,rawat2025llmmas,pan2026graphdtgpt,zhang2025llmagentschema}，不是无监督飞控或静默篡改孪生的许可证。没有检索的 ReAct 循环只是多几步的参数叙事~\cite{yao2023react}；没有孪生记忆的工具调用可能打错宿主；不受控的记忆流可能把推测积成``事实''~\cite{park2023generativeagents}；因此写入门禁应落在 harness 里~\cite{luo2025llmdtrobot,gao2026ifcagent}。

\begin{table}[htbp]
\centering
\caption{专门化到立面巡检的记忆阶梯。}
\label{tab:ch8-memory}
\begin{tabular}{llp{5.2cm}}
\toprule
层 & 保存什么 & 巡检用途 \\
\midrule
参数记忆 & 模型通用知识 & 语言与常见术语 \\
上下文记忆 & 提示 + 对话 & 当前问题状态 \\
检索记忆（RAG） & 手册、观测卡片 & 带版本文本证据 \\
孪生库 & Asset-ID + 拓扑 + 时间 & 权威现场记忆 \\
\bottomrule
\end{tabular}
\end{table}

\begin{figure}[htbp]
\centering
\begin{tikzpicture}[node distance=6mm and 7mm]
  \node[ch8boxalt, minimum width=2.5cm] (p) {参数记忆};
  \node[ch8box, minimum width=2.5cm, right=of p] (c) {上下文记忆};
  \node[ch8boxrag, minimum width=2.5cm, right=of c] (r) {检索记忆\\(RAG)};
  \node[ch8boxok, minimum width=2.5cm, right=of r] (t) {孪生库\\(Asset-ID)};
  \draw[ch8arr] (p) -- (c);
  \draw[ch8arr] (c) -- (r);
  \draw[ch8arr] (r) -- (t);
  \node[ch8tiny, below=1mm of p] {通用权重};
  \node[ch8tiny, below=1mm of c] {会话窗口};
  \node[ch8tiny, below=1mm of r] {带版本文档};
  \node[ch8tiny, below=1mm of t] {权威现场记忆};
\end{tikzpicture}
\caption{立面巡检的记忆阶梯：孪生库是权威；参数记忆永远不能当作现场数据库。}
\label{fig:ch8-memory-ladder}
\end{figure}

智能体记忆研究常用``生成、更新、演进''闭环描述经验条目~\cite{park2023generativeagents,yao2023react,wang2024surveyagents}。对立面巡检，我们保留同一符号，但强制\emph{外部化}演进：变的是孪生与 harness，模型权重保持冻结。令 $S_t$ 为查询~$t$ 的推理态，涵盖意图、工具 TRACE 与校验结果；$M_t$ 为外部记忆，涵盖孪生库、版本钉扎与过程轨迹。则
\begin{equation}
\label{eq:ch8-mem-loop}
M_{t+1}=\mathrm{Update}\!\bigl(\mathrm{Generate}(S_t,M_t),\,M_t\bigr),\qquad
\pi_{t+1}=\mathrm{Evolve}_{\mathrm{ext}}(\pi_t,M_{t+1}),
\end{equation}
其中 $\mathrm{Generate}$ 可产出护照摘要、驳回草稿备注或复访提示；$\mathrm{Update}$ 仅经有门写回或过程日志合并；$\mathrm{Evolve}_{\mathrm{ext}}$ 改变 harness 内的工具策略、版本钉扎或检索包，并禁止在现场文本上静默微调~\cite{lewis2020,luo2025llmdtrobot,gao2026ifcagent}。实践上，会话别名保持短期状态，带模态标签的观测形成可检索卡片，资产护照提供语义摘要，批准后的轨迹保存过程记忆；孪生库、版本钉扎与门禁共同更新运行包，同时把现场事实留在模型权重之外。

这三层意在组合，而非互斥。再看：「西立面中层空鼓是否与上层渍迹相关，本季度要不要安排吊篮？」L1 解析混合问句并强制模板；L2 工具取种子护照并扩展 BIM 邻域，RAG 装入卡片与钉扎手册条款，结构化生成输出 ID 属于证据集的草稿；L3 保证读的是长期孪生而非昨日聊天，且工单或复访提议须经人工门。跳过 L2，答案沦为参数讲故事；跳过结构校验与门禁，流畅文本会漏进 CMMS。这一组合正是「生成式 AI 决策支持」在立面运维上的可操作读法~\cite{lewis2020,yao2023react,ji2023}。

把三层真正黏成可上线系统的，是\emph{模型 harness}，即运行时驾驭层。在软件测试里，harness 指在受控条件下驱动被测组件的脚手架；在当代 LLM 工程里，同一词越来越多地指包裹可替换语言模型的、由代码拥有的控制层：它组装提示、选择检索单元、暴露类型化工具、按契约校验输出、记录轨迹，并决定何时拒答或升级。智能体综述已把规划、记忆与工具视为自主性的要素~\cite{wang2024surveyagents,li2024surveymultiagent}；推理，行动循环与工具学习则说明模型可以请求动作，却不必拥有世界状态~\cite{yao2023react,schick2023toolformer}。HELM 一类整体评价框架进一步提醒读者：场景、指标与标准化运行条件本身就是比较用的 harness。单一准确率数字并不足够~\cite{liang2023helm}。对立面巡检，这一思想可专门化为：Asset-ID 契约、版本钉扎、模态标签与写入门禁应住在 harness 里，这样更换底层 LLM 不会静默改写运维政策。

巡检 harness 可拆成五个协同部分，见表~\ref{tab:ch8-harness}。\emph{接入与路由}负责意图分类，并在生成前把自然语言位置解析到 GeoBIM 别名。\emph{证据装配}加载护照、邻域、观测卡片与钉扎手册，即第 L2 层的 RAG 与打包职责~\cite{lewis2020,peng2024,krutli2025fmrags}。\emph{工具面}只暴露少量类型化动作：读取护照、扩展邻域、列出观测、起草报告、提议复访，并拒绝自由 SQL 或飞控指令；面向 IFC、能耗与运维机器人的 AEC 工具智能体同样把推理与世界动作分开并置于批准之下~\cite{gao2026ifcagent,zhang2025llmagentschema,luo2025llmdtrobot,guo2025bimllmquery}。\emph{契约校验器}在起草后检查可机检义务：点名的 Asset-ID 必须出现在检索包中；数值缺陷主张必须指向观测字段；模态冲突必须披露；写入或调度必须经人工门~\cite{ji2023,pan2024b,buga2025llmfmsql}。\emph{轨迹与评价}保存跑过哪些工具、钉扎了哪一版手册、哪条校验通过或失败，使答案可重放、可评分，这是对多指标、多场景评价的巡检读法。聊天流畅度并不是充分指标~\cite{liang2023helm,shuster2021}。自适应检索重试应作为 harness 内的受控循环。临时即兴改提示应被禁止~\cite{asai2024selfrag,jeong2024adaptive}。

\begin{table}[htbp]
\centering
\caption{巡检模型 harness：包裹可替换 LLM 的代码侧部件。}
\label{tab:ch8-harness}
\begin{tabular}{llp{4.8cm}}
\toprule
部件 & 负责什么 & 巡检例子 \\
\midrule
接入/路由 & 意图、Asset-ID 解析 & 西侧跨 $\rightarrow$ GeoBIM 种子 \\
证据装配 & RAG + 打包 & 护照、邻域、手册 \\
工具面 & 类型化 API + 门禁 & 扩展邻域；禁飞控指令 \\
契约校验 & 引用或拒答检查 & ID $\subseteq$ 检索包；模态标记 \\
轨迹/评价 & 可重放日志 + 指标 & 工具 TRACE；类 HELM 场景 \\
\bottomrule
\end{tabular}
\end{table}

\begin{figure}[htbp]
\centering
\begin{tikzpicture}[yscale=0.90]
  \def\ch8R{3.35cm}
  \node[font=\scriptsize\itshape, text=ch8mute, align=center] at (0,0)
    {代码拥有的\\查询环};
  \node[ch8box, minimum width=2.0cm, minimum height=0.88cm, font=\footnotesize,
        inner sep=3pt] (in) at (90:\ch8R) {接入/\\路由};
  \node[ch8boxrag, minimum width=2.0cm, minimum height=0.88cm, font=\footnotesize,
        inner sep=3pt] (ev) at (30:\ch8R) {证据\\装配};
  \node[ch8boxok, minimum width=2.0cm, minimum height=0.88cm, font=\footnotesize,
        inner sep=3pt] (tool) at (-30:\ch8R) {工具面};
  \node[ch8boxllm, minimum width=2.0cm, minimum height=0.88cm, font=\footnotesize,
        inner sep=3pt] (llm) at (-90:\ch8R) {可替换\\LLM};
  \node[ch8boxgate, minimum width=2.0cm, minimum height=0.88cm, font=\footnotesize,
        inner sep=3pt] (val) at (-150:\ch8R) {契约\\校验};
  \node[ch8boxalt, minimum width=2.0cm, minimum height=0.88cm, font=\footnotesize,
        inner sep=3pt] (tr) at (150:\ch8R) {轨迹/\\评价};
  \foreach \a/\b in {in/ev, ev/tool, tool/llm, llm/val, val/tr, tr/in}
    \draw[ch8arr, shorten >=1pt, shorten <=1pt] (\a) -- (\b);
  \foreach \name/\n in {in/1, ev/2, tool/3, llm/4, val/5, tr/6}
    \node[ch8step] at ($(\name.north west)+(-0.08,0.08)$) {\n};
\end{tikzpicture}
\caption{巡检模型 harness 的六步查询环。步骤 1--3 与 5--6 由代码拥有；步骤 4 为可替换 LLM。每个设施问题都走完这一环。}
\label{fig:ch8-harness-fig}
\end{figure}

按顺时针读图~\ref{fig:ch8-harness-fig}：接入先解析意图与 Asset-ID 种子；证据装配打包护照、邻域与钉扎手册；工具面只允许类型化的孪生读取与起草动作；可替换 LLM 在该证据包内组织语言；契约校验强制引用或拒答与模态披露；轨迹与评价使本次运行可重放，下一查询仍走同一门禁。虚线 LLM 框是故意的：权重可换，外围契约不宜随换。

由此得到两条采购与教学推论。其一，纯提示原型还不是 harness：``务必引用 Asset-ID''这类指令可被流畅模型忽略，而``ID 子集检查失败则禁止交付''的校验器才是承重结构。其二，模型替换应是 harness 测试。模型替换不应变成项目重写，若校验器、工具 schema 与孪生接口固定，更换托管 LLM 应只改变措辞，而不改变哪些来源合格、写入门是否存在~\cite{liang2023helm,saka2024gptconstruction}。业界关于 harness engineering 的讨论强调把行为从提示迁到版本化契约、清单与轨迹；本章保留这一直觉，但只引用已发表的接地、工具、智能体与评价结果，并把每条契约映射到立面运维产物，而非通用聊天演示。简言之：L1--L3 命名能力；harness 是在每次查询上强制执行这些能力的软件。

近年面向建筑的 LLM 工作同样把重心从单纯模型选型转向系统设计。数据就绪、信任、互操作与人机监督现已置于参数规模之前~\cite{ma2026tenquestions}：立面运维成败取决于 Asset-ID、带版本手册与批准门是否就绪。已发表的图孪生、BIM 检索、IFC 工具与运维机器人智能体，在邻近场景中展示同一习惯：把类型化结构与门禁放在权重之外，让 LLM 作为 harness 内可替换的组合器~\cite{pan2026graphdtgpt,lu2026highwaymas,gao2026ifcagent,luo2025llmdtrobot,hellin2025bimmas,quan2026bimfacrag}。这些系统都不是开箱即用的仓立面巡检器；可迁移的是控制模式。

幻觉综述解释了缺乏纪律时会发生什么：流畅系统会虚构实体、数字与引用~\cite{ji2023}。在巡检中，高代价失效非常具体，虚构 Asset-ID、缺乏可归因证据的维修建议、抹掉 RGB 与红外观测分歧的摘要。因此接地远不止``写得好的提示''。约束 schema、检索及其自适应变体、工具型推理，行动循环、类型化结构化知识在 NLP 与 AEC 中反复出现~\cite{lewis2020,gao2023hyde,asai2024selfrag,jeong2024adaptive,peng2024,yao2023react,schick2023toolformer,pan2024b}。面向图与层次的检索提醒读者：文档相似不等于建筑邻接~\cite{peng2024,sarthi2024raptor,sun2024tog}。语义嵌入会把``潮湿、渍迹、渗水''聚在一起；建筑空间则按接缝、楼层与排水路径聚集板块。混合检索可以在\emph{已被拓扑约束的候选集}内用嵌入排序症状摘要，却不得用嵌入为传播问题\emph{发明}候选集~\cite{peng2024,jeon2025hybridrag,khattab2020}。当系统规则、检索手册、孪生事实与用户请求冲突时，优先级应明确：安全与写入门禁优先；孪生事实与观测字段其次；钉扎手册版本再次；用户文风偏好最后~\cite{saka2024gptconstruction}。

土木与建筑应用已经在使用这些模式。BIM 助手与自然语言模型查询改善导航；规范、安全、缺陷与设施 RAG 改善版本忠实；schema、类型化工具与批准使语言接口可运维~\cite{zheng2023bimassistant,joffe2025buildingcodes,lee2024ragsafety,jeon2025hybridrag,krutli2025fmrags,gao2026ifcagent,buga2025llmfmsql}。高分辨率立面数据与物理量化仍决定语言可诚实宣称的粒度~\cite{zhao2024facade,zhao2026physical}。对本书而言，第2--7章保有飞行、几何、观测与 GeoBIM 身份；第8章把这些产物组织成证据包~\cite{zhang2025,yang2025detreconreg,han2025}。因此默认应在资产护照与 BIM 邻域上检索，手册按版本钉扎；公开规范题上的成功，不能证明本地立面手册也可靠~\cite{lee2024ragsafety,ma2026tenquestions}。

记忆失效把阶梯讲具体。若经理澄清``西侧''指装卸台立面而次日对话忘记，修复是持久化别名或强制 GeoBIM 种子选择。团队不应指望模型自行``记得''。若 RAG 返回过期维护 PDF，修复是版本钉扎与语料冒烟测试~\cite{krutli2025fmrags,joffe2025buildingcodes}。若模型凭通用网络知识发明缝长，修复是禁止数值缺陷主张，除非观测字段提供~\cite{ji2023,zhao2026physical}。若孪生有正确 Asset-ID 但编排器根本没取，那是基础设施故障，不是文笔问题~\cite{lewis2020,peng2024}。LLM 起草后谁可写入孪生、对话澄清保留多久、两份 PDF 冲突时哪一版手册权威，这些都是基础问题，不是事后 IT 政策。对立面运维，建议简单规则：孪生库写入有门；对话窗口易失；RAG 索引带版本；参数记忆永不当作现场数据库~\cite{luo2025llmdtrobot,gao2026ifcagent,rawat2025llmmas}。

选择巡检用的 LLM 技术栈，同样更是 harness 问题而非单纯模型选型问题。团队应评估上文五个 harness 部件，尤其是校验器与工具门禁是否由代码拥有，以及原生工具调用可靠性、护照集合上的长上下文行为、约束解码下的拒答意愿、建筑数据的本地/云约束，以及每个接地答案的成本而非闲聊 token 成本~\cite{saka2024gptconstruction,wang2024surveyagents,liang2023helm}。放在严格引用或拒答 harness 里的较小模型，常常胜过在无校验闲聊环里即兴编造 Asset-ID 的较大模型。即便在完整实验计分之前，手检十个答案已能区分接地作业与聊天演示：列出的 Asset-ID $\subseteq$ 检索集；假说已标记；模态标签保留；缺邻域时拒答~\cite{ji2023,shuster2021,lewis2020}。教学结论因此是三项设计选择：在可维护资产及其邻域上检索；强制引用检索集合中出现的 ID；把不可逆动作放在批准门之后。第\ref{sec:ch8-dt-to-language}--\ref{sec:ch8-multimodal}节接着说明孪生记录与多模态证据如何让 L2 与 L3 落地。



概括基础部分：第 L1--L3 层命名现成模型在控制下能做什么，harness 命名在每次查询上强制这些行为的软件。跳过校验器与工具门禁的读者，无论模型多大，都会重新掉进参数叙事~\cite{ji2023,liang2023helm,ma2026tenquestions}。本章把语言层写成可核对、可交接的决策接口，而不是一段流畅但不可审计的聊天演示。读罢全文，读者应能说明：孪生若缺表~\ref{tab:ch8-twin-contract} 中某项最低前提，答案会在何处失真；一个立面问题该归入 L1--L3 的哪一层，并由表~\ref{tab:ch8-harness} 中哪一段 harness 承托；传播题与历史记录题如何用护照与 BIM 邻域拼成``带引用或拒答''的证据包；第~\ref{sec:ch8-multimodal}节里 VLM 与第~6章检测器各管哪一段；以及 twin 写入或复访提议怎样在人工门禁前留下一条可回放的工具 TRACE~\cite{peng2024,yao2023react,liang2023helm,hellin2025bimmas}。表~\ref{tab:ch8-outcomes} 将上述能力收成核对清单。它们界定的是跨章决策接口：第~2--7章供给几何、检测与孪生物料，本章说明语言侧如何装配与门禁，并不再引入新的检测器。

\begin{table}[htbp]
\centering
\caption{第8章可核对能力清单。}
\label{tab:ch8-outcomes}
\begin{tabular}{llp{5.0cm}}
\toprule
核对项 & 应掌握的内容 & 合格判据 \\
\midrule
孪生契约 & 表~\ref{tab:ch8-twin-contract} 五项前提 & 每项对应一种缺失失效 \\
L1--L3 + harness & 问题所落层级与 harness 部件 & 能用表~\ref{tab:ch8-harness} 的语言表述 \\
证据装配 & 传播题或历史记录题的证据包 & ID$\subseteq$包，或明确拒答 \\
多模态分工 & VLM 与检测器的边界 & 检测器写观测，VLM 只读裁剪 \\
有界行动 & 带门禁的工具 TRACE & 无静默写孪生 \\
\bottomrule
\end{tabular}
\end{table}

下一节把这些层落实为孪生记录：观测卡片、资产护照与按意图装配。

\section{从数字孪生到语言推理}
\label{sec:ch8-dt-to-language}

在巡检孪生上做语言推理，不是自由叙事，而是装配前几章已经写好的证据，这也是第~\ref{sec:ch8-llm-foundations}节第 L2--L3 层的操作底座。第~7章已在 GeoBIM 层给出可维护宿主与设施台账身份；本章把这类稳定编码简称为 \emph{Asset-ID}，指孪生中可长期引用、可与工单和 CMMS 对齐的建筑对象标识，而不是临时文件名或检测框标签。只有在对象身份、出处、邻域、时间与模态都能被机器检查时，语言层才谈得上``接地''。

\begin{table}[htbp]
\centering
\caption{在 LLM 支撑立面运维之前，孪生至少应具备的承诺。}
\label{tab:ch8-twin-contract}
\begin{tabular}{llp{5.4cm}}
\toprule
承诺 & 典型载体 & 缺失时的失效 \\
\midrule
稳定资产标识（Asset-ID） & BIM / GeoBIM 宿主 & 虚构板块名称 \\
观测出处 & 检测日志 + 影像路径 & 不可审计主张 \\
拓扑 / 邻域 & 立面邻接 & 无法传播阅读 \\
时间戳 & 航次元数据 & 伪``当前状态''叙述 \\
模态标签 & RGB / IR / \ldots & 静默揉合证据 \\
\bottomrule
\end{tabular}
\end{table}

表~\ref{tab:ch8-twin-contract} 不是 IT 心愿单，而是现场手艺的数字化延续。资深巡检员先拍照再写字，先点名跨位再诊断，拒绝签署无法到现场核对的报告；决策可用孪生只是让这些习惯在跨航次尺度上可被机器检查~\cite{zhang2022,chen2023facility}。地理尺度把同一点说得更尖：校园或住宅区会放大标识冲突与手册版本冲突，城市尺度身份应在 GIS 层保持一致，而不应事后用巧妙散文``修复''~\cite{zhang2025,han2025}。

即便实现各不相同，概念上仍宜分清三层。\emph{观测记录}是带传感出处的一次缺陷目击，涵盖影像、证据、语义与宿主挂接，教学上常打包为集成缺陷画像；\emph{资产护照}聚合某一可维护宿主上的观测与几何/拓扑，作为供检索的集成资产画像；\emph{查询装配}决定生成前哪些护照与邻域进入上下文，沿立面邻接扩展邻域则支持竖向传播阅读。感知写观测，运维按资产思考，语言模型不得发明二者~\cite{pan2024b,zheng2023bimassistant}。

在无人机巡检实践中，有用的观测通常包含宿主 Asset-ID、时间、模态，如 RGB、红外等、检测器签名、位姿或表面坐标，以及影像路径。任何 LLM 介入之前，配准就应把检测挂到 BIM 宿主上~\cite{zhang2025,yang2025detreconreg}。自然语言位置短语帮助模型，例如``西立面、装卸台上方第三跨''，度量坐标留给现场导航。资产护照再附上症状摘要、用于严重度归一化的几何，以及上/下/侧邻接。这些链接使传播式阅读成为可能，例如把中层空鼓与上层潮湿一并考虑，而不把文档相似图误当作结构图或 BIM 邻接图~\cite{peng2024}。

\begin{figure}[htbp]
\centering
\includegraphics[width=\linewidth]{idp_construction_pipeline.pdf}
\caption{从航次影像与检测器输出构建语言可用巡检记录：带模态与证据指针的类型化观测挂接到可维护宿主，再聚合为检索可加载的资产中心护照。各实现命名不同；教学模式仍是观测$\rightarrow$宿主挂接$\rightarrow$护照。}
\label{fig:ch8-obs-passport}
\end{figure}

图~\ref{fig:ch8-obs-passport}区分感知所写的一次带出处目击，与设施语言所需的含邻域与近期观测的宿主级护照。第4章提供几何与 BIM 宿主作为挂接点，第5--6章提供带标签观测与置信度，第7章提供城市尺度 Asset-ID；本章消费这些产物，使检索、记忆、报告和行动都可追溯。

面向语言使用的最小观测 schema 包括观测 ID、宿主 Asset-ID、航次或采集时间、模态标签、第6章给出的缺陷类型与置信度、可选的由几何导出的严重度特征、影像或裁剪 URI，以及简短可读位置短语。可选字段包括天气标记、巡检员备注，以及同一宿主上历史观测链接。绝不能只靠自由文本维系的是宿主 Asset-ID。采集阶段临时起的影像文件名或文件夹标签，记录的是文件如何存放，并不是设施台账编码。若影像到宿主的唯一链接只是这类文件名约定，检索终将发明邻居。企业 RAG 常常切分 PDF，但巡检问题通常关于可维护宿主及其邻域。因此护照应为每个 Asset-ID 打包紧凑症状摘要、近期观测 ID、用于严重度归一化的几何摘要，以及显式邻域列表。严重度归一化很重要，因为小板块上 20\,cm 裂缝与大墙跨上 20\,cm 裂缝并不是同一运维对象；语言系统应从孪生几何继承该归一化。语言层不得自行发明该归一化~\cite{zhao2026physical}。邻域列表应来自 BIM 或 GeoBIM 邻接。报告句子之间的嵌入相似不能充当邻域。

查询装配是巡检不同于通用企业 RAG 之处。先把已构建记录写清楚，再让意图选择检索包。宿主~$a$ 的资产护照打包身份、几何摘要~$G(a)$、一跳 BIM 邻域~$N_1(a)$、观测集~$O(a)$与紧凑症状摘要~$d(a)$：
\begin{equation}
\label{eq:ch8-passport}
P(a)=\bigl\langle a,\,G(a),\,N_1(a),\,O(a),\,d(a)\bigr\rangle.
\end{equation}
由问题~$q$ 解析得到的种子 Asset-ID 集合~$A_q$，以及由意图~$\iota$ 决定的跳数~$k(\iota)$：局部查找为 0，传播为 1，BIM 邻域为
\begin{equation}
\label{eq:ch8-khop}
\mathcal{N}_{k}(A_q)=\bigl\{v:\exists\,a\in A_q,\;\mathrm{dist}_{\mathrm{BIM}}(a,v)\le k\bigr\},
\end{equation}
供生成使用的证据包为
\begin{equation}
\label{eq:ch8-package}
\mathcal{C}(q)=\mathrm{Pack}\!\Bigl(\bigl\{P(v):v\in\mathcal{N}_{k(\iota)}(A_q)\bigr\},\,R_{\mathrm{man}}(q),\,e\Bigr),
\end{equation}
其中 $R_{\mathrm{man}}(q)$ 为可选的带版本手册切片，$e$ 为版本钉扎本身~\cite{lewis2020,peng2024,joffe2025buildingcodes}。局部查找，例如``资产~$A$ 上月登记了哪些缺陷？''取 $k=0$；传播类问题取沿立面拓扑的邻域，而不是仅共享词汇的 top-$k$ 文本块；楼层或立面审计则用第7章空间过滤替代跳数扩展。因此意图改变的是进入~$\mathcal{C}(q)$ 的\emph{哪些}孪生记录，而不只是提示措辞~\cite{lewis2020,peng2024}。当问题把现场证据与规范语言混在一起时，$R_{\mathrm{man}}$ 可并入带版本手册与安全条款~\cite{joffe2025buildingcodes,lee2024ragsafety,krutli2025fmrags}。

\begin{table}[htbp]
\centering
\caption{常见查询意图及其应加载的孪生记录。}
\label{tab:ch8-intents}
\begin{tabular}{llp{4.8cm}}
\toprule
意图 & 主要记录 & 检索单元错误时的典型失效 \\
\midrule
局部查找 & passport($A$) & 漏掉近期观测 \\
资产健康 & passport($A$) + 历史 & 空泛``立面良好'' \\
传播阅读 & $A$ + 邻域 & PDF 块假关联 \\
立面审计 & 空间过滤集合 & 挑拣轶事 \\
规范 + 现场 & 手册 + 护照 & 有条款无证据 \\
\bottomrule
\end{tabular}
\end{table}

一种反复出现的实现模式把观测组织为资产护照，并沿 BIM 邻接扩展邻域以支持多跳诊断。已发表的同类工作在其他语料上共享接地直觉：缺陷混合 RAG~\cite{jeon2025hybridrag}、规范与安全 RAG~\cite{joffe2025buildingcodes,lee2024ragsafety}、设施文档 RAG~\cite{krutli2025fmrags}、BIM 助手~\cite{zheng2023bimassistant}。这里强调的是模式：巡检问题以资产为中心，检索单元就应以资产为中心。当装配无法支撑某主张时，正确行为是拒答或升级给人工巡检员。自信的空话应被拒答机制拦住。

微型示例。假设经理问：``西侧竖向栈的渗水是否可能由上方传播？''装配应加载目标护照、上方与侧向邻域、带 RGB/红外标签的观测卡片，以及挂到这些 Asset-ID 的潮湿相关维护说明，如有则一并加载。生成仅在邻域观测确实存在时才可提出路径假说；否则必须说明孪生尚不支持传播判断。这才是``在孪生上做语言推理''。``在孪生旁边讲故事''不符合本章要求。

实现团队常问护照该有多大。太小，传播阅读会饿死；太大，上下文会被整面立面的无关邻域填满。实用教学启发式是按意图条件扩展：局部查找零跳；传播沿 BIM 邻接一跳或小 $k$；立面审计用第7章空间过滤而非无限跳~\cite{peng2024,zhang2025}。仅当问题提到合规、安全或维修标准时，才检索手册条款~\cite{joffe2025buildingcodes,lee2024ragsafety}。

索引选择应服从同一启发式。护照摘要上的向量索引支持模糊症状搜索，例如``空鼓且上方有渍迹的板块''。Asset-ID 上的图或关系索引支持精确邻域扩展。观测 ID 上的关键词索引支持审计查找。混合检索之所以有吸引力，正因为巡检问题把模糊语言与精确标识混在一起~\cite{jeon2025hybridrag,krutli2025fmrags}。但拓扑权威性不能被混合掉：当问题关于传播时，嵌入相似不得静默取代 BIM 邻接。

拒答应有同等教学时间。工程师在现场本来就会说``证据不足''；语言系统必须学会同一习惯。令 $\hat{y}$ 为草稿答案，$\mathrm{IDs}(\cdot)$ 为字符串或检索包中点名的 Asset-ID 与观测 ID。可机检的引用或拒答门禁为
\begin{equation}
\label{eq:ch8-cite}
\mathrm{IDs}(\hat{y})\subseteq\mathrm{IDs}\bigl(\mathcal{C}(q)\bigr)
\quad\text{否则返回 }\texttt{ABSTAIN}.
\end{equation}
严重度主张还必须指向 $\mathcal{C}(q)$ 内 $O(\cdot)$ 的成员；任何传播路径必须是式~\eqref{eq:ch8-khop} 已扩展 BIM 邻域中的一条通路。若检查失败，系统返回结构化拒绝，并说明缺失证据类型，例如无上层邻居、模态冲突或航次过旧。这比一段无法在孪生界面打开的流畅段落更有价值~\cite{ji2023,shuster2021}。

第二个微型示例澄清立面审计。经理问：``汇总三层以上、需要吊篮接近的西立面问题。''装配应应用第7章空间过滤，加载所得护照集合，保留模态标签，并按几何继承的严重度特征排序。它不应叙述整栋楼，也不应丢掉 Asset-ID 只剩楼层散文。输出是排程用的决策交付物，不是游记。



经理常问某渍迹是否``新的''。查询历史维修记录是决策层的一等能力，也是一等装配意图。跨航次比较需要跨时间同一 Asset-ID、可比的模态覆盖、可用时的既往观测 ID 或工单链接，以及对缺失既往记录的显式处理。若上一航次无红外，系统不得发明热差值；若几何重网格后 Asset-ID 被随意重映射，语言就会比较错宿主。自然语言问题如``2024 航次之后西立面修过什么？''应检索带版本孪生记忆与任何已批准报告链接；缺 ID 时必须拒答，而不得编造维修叙事~\cite{zhang2025,yang2025detreconreg,chen2023facility}。纵向纪律因此首先是孪生卫生，其次才是模型聪明~\cite{zhang2025,yang2025detreconreg,zhao2026physical}。感知可能输出置信度；几何可能支持长度或面积特征；运维可能需要序数严重度以便排程。语言层不应把序数分诊写成虚假的毫米级确定性，更宜使用绑定字段的表述，例如``severity=medium，来自几何归一化裂缝长度；confidence=0.71，来自检测器''。电影式形容词应被避免。物理量化管线说明数字为何需要出处~\cite{zhao2024facade,zhao2026physical}；LLM 的职责是把出处带进句子。语言层的职责也不是装饰出处。

算法~\ref{alg:ch8-assembly} 把装配写成带编号的步骤。它故意无聊：每个分支都是学生想耍聪明的地方，也是巡检系统通常折断的地方~\cite{lewis2020,peng2024,jeon2025hybridrag}。

\begin{algorithm}[htbp]
\caption{面向接地巡检答案的意图条件孪生装配。}
\label{alg:ch8-assembly}
\begin{algorithmic}[1]
\Require 自然语言查询 $q$
\Ensure 草稿答案，或结构化拒答码
\State $\mathit{intent} \gets \textsc{Classify}(q)$
\State $\mathit{seeds} \gets \textsc{ResolveAssets}(q)$ \Comment{否则 \texttt{ABSTAIN:asset\_unresolved}}
\State $\mathit{packs} \gets \{\textsc{Passport}(s): s\in\mathit{seeds}\}$
\If{$\mathit{intent}=\texttt{propagation}$}
  \State $\mathit{packs} \gets \mathit{packs}\cup\textsc{Expand}(\mathit{seeds},k{=}1)$
\EndIf
\If{$\mathit{intent}=\texttt{elevation\_audit}$}
  \State $\mathit{packs} \gets \textsc{SpatialFilter}(\mathit{packs},q)$
\EndIf
\State $\mathit{obs} \gets \textsc{ListObservations}(\mathit{packs})$
\State 若 $\textsc{NeedsCode}(q)$ 则 $\mathit{manuals} \gets \textsc{RetrieveManuals}(q)$，否则 $\mathit{manuals} \gets \emptyset$
\State $\mathcal{C}(q) \gets \{\mathit{intent},\mathit{packs},\mathit{obs},\mathit{manuals},\mathit{edition\_pin}\}$
\State $\hat{y} \gets \textsc{Generate}(\mathcal{C}(q))$ \Comment{引用或拒答}
\State $\textsc{ValidateIDs}(\hat{y},\mathcal{C}(q))$ \Comment{否则 \texttt{ABSTAIN:id\_not\_in\_package}}
\State \Return $\hat{y}$
\end{algorithmic}
\end{algorithm}

自然语言位置仍然含糊：``第三跨''在重网格后可能对应不同宿主。必须先对 GeoBIM 别名解析，再取护照。模糊匹配可以提出候选；当两个候选分数接近时，不得静默选一。消歧提问是成功的系统行为，不是 UX 失败~\cite{zhang2025,han2025,guo2025bimllmquery}。护照回答``这个宿主此刻被怎样概括''；观测回答``哪一次目击、何种模态、何种置信度''。生成若只引护照不引观测，审计员无法打开影像；若只引观测不引护照，排程员无法落到可维护对象。双引是默认。这是默认纪律~\cite{pan2024b,zhao2026physical}。立面审计用第7章空间过滤得到宿主集合；传播用邻接扩展得到路径。先过滤再对集合内每个宿主做传播，是允许的组合；用嵌入``找相似板块''冒充过滤或扩展，则不允许~\cite{peng2024,zhang2025}。当结构化库缺少路径时拒绝用无约束文本硬答，这一纪律也出现在近年孪生查询系统中~\cite{pan2026graphdtgpt,buga2025llmfmsql}：对立面巡检，``孪生中无上层邻域''比流畅渗水故事更是好产品行为。

2024--2025 年已发表检索实践也澄清了不必重造什么。混合缺陷管理 RAG 与安全手册 RAG 表明：语料带版本、运维语言嘈杂时仍有收益~\cite{jeon2025hybridrag,lee2024ragsafety}；规范与法规混合检索表明：条款检索仍须会合现场证据，建议才能进入运维~\cite{joffe2025buildingcodes,he2025regulationhybrid,yang2024promptcode}；设施文档 RAG 与结构化 FM 查询表明：只有存在 schema 与 Asset-ID 链接时，建筑语料才能规模化~\cite{krutli2025fmrags,buga2025llmfmsql,arslan2024ragenergy}。把这些结果迁到无人机立面运维，意味着以观测卡片与资产护照为一等检索单元，手册为次级钉扎上下文，并把拓扑扩展当作一等装配步骤。拓扑扩展也不是嵌入的事后补丁~\cite{peng2024,pan2024b,ma2026tenquestions}。




把装配纪律再展开一层教学叙述，有助于设施与软件两边对齐词汇。局部查找题只应装入目标护照与近期观测；若系统顺手塞进整面立面的护照，上下文会被噪声占满，引用或拒答校验也会变得难以解释~\cite{lewis2020,peng2024}。传播题必须先有拓扑边，再谈症状相似；若边缺失，正确产品行为是拒答并点名``无邻居''。用嵌入把对面立面的``渍迹''拉进来应被禁止~\cite{peng2024,jeon2025hybridrag,pan2026graphdtgpt}。立面审计题应走第7章空间过滤，输出可调度的 Asset-ID 列表。旅行游记式散文不符合交付要求~\cite{zhang2025,buga2025llmfmsql}。规范+现场题必须同时钉扎手册版本与现场护照；只有条款没有观测，或只有观测没有条款，都应在答案中显式披露，而不能用流畅句子掩盖缺口~\cite{joffe2025buildingcodes,lee2024ragsafety,he2025regulationhybrid,quan2026bimfacrag}。跨航次比较题则把身份卫生推到前台：Asset-ID 重映射失败时，``是否更严重''必须拒答，这比给出一个漂亮差值更负责任~\cite{zhang2025,yang2025detreconreg,zhao2026physical}。

算法~\ref{alg:ch8-assembly} 故意写得``无趣''：每条分支都是实现上常想走捷径、系统却最容易崩的地方。意图分类错了，会装错包；种子解析错了，会点名错误宿主；扩展跳数过大，会把整栋楼塞进提示；手册未钉扎，会把两版条款折中成四不像。校验器放在生成之后，是为了抓住模型仍然可能滑出的虚构 ID，提示词里的``请引用''从来不是门禁~\cite{ji2023,liang2023helm,ma2026tenquestions}。BIM 自然语言多智能体检索与精简检索智能体表明：导航模型对象可以很聪明，但巡检还要在导航之后做拓扑诊断与模态披露~\cite{hellin2025bimmas,koh2026bimllmretr,liu2025bimcoder}。设施 RAG 表明：文档库再大，若不回链 Asset-ID，就仍难进入 CMMS~\cite{krutli2025fmrags,quan2026bimfacrag,arslan2024ragenergy}。把这些文献读成装配策略，装配层才会稳定。




表~\ref{tab:ch8-assembly-fail} 把同一纪律压成课堂与采购语言。每一行把常见捷径与 harness 应发出的结构化拒答配对。学生可用冻结孪生上``系统能否产出这些拒答''给演示打分；实现者可把右列变成算法~\ref{alg:ch8-assembly} 旁的单元测试~\cite{ji2023,peng2024,liang2023helm}。

\begin{table}[htbp]
\centering
\caption{常见装配捷径与应阻止它们的拒答。}
\label{tab:ch8-assembly-fail}
\begin{tabular}{llp{4.4cm}}
\toprule
捷径 & 会坏什么 & 优选拒答 \\
\midrule
模糊宿主静默选定 & 错误 Asset-ID 工单 & \texttt{ABSTAIN:asset\_ambiguous} \\
嵌入``相似渍迹'' & 假渗水路径 & \texttt{ABSTAIN:no\_neighbour} \\
整面立面倾倒进上下文 & 引用无法核对 & 收缩为过滤后集合 \\
手册未钉扎 & 版本折中 & \texttt{ABSTAIN:edition\_unpinned} \\
跳过 ID 校验 & 散文虚构板块 & \texttt{ABSTAIN:id\_not\_in\_package} \\
\bottomrule
\end{tabular}
\end{table}

概括孪生到语言的装配：意图决定装哪些护照与邻域；引用或拒答决定草稿能否离开 harness；GeoBIM 解析决定哪个宿主有资格被点名~\cite{lewis2020,peng2024,zhang2025}。表~\ref{tab:ch8-assembly-fail} 使这些决策可核对；下一节再在同一契约下加入模态标签与有界智能体。

\section{多模态融合}
\label{sec:ch8-multimodal}

立面项目天然多模态：航次常同时产生 RGB 影像、红外通道与文本元数据。在本章中，\emph{多模态融合}指的是：把第6章感知已经写出、并由第4与第7章完成挂接的带模态标签观测记录，纪律地装配进接地 LLM 检索与生成。缺陷发现仍是第6章感知职责；语言负责在引用规则下检索、解释、拒答与行动提议。

视觉—语言模型（VLM）属于全书决策支持阶段的一部分，但并不取代第5--6章的缺陷检测器。VLM 可以帮助巡检员叙述已经挂到 Asset-ID 上、可打开的裁剪，把自由文本备注与检索到的护照对齐，或起草仍引用观测 URI 的图注。它不得发明宿主编码，不得静默改写检测器类别，也不得把聊天截图当作新的孪生事实。因此本章讲授的分工很简单：第6章写类型化观测；VLM 若使用，则是在与 LLM 同一套引用或拒答门禁下阅读这些裁剪的可选阅读器；孪生仍是权威库~\cite{zhao2024facade,zhao2026physical,ji2023}。这一分工把视觉缺陷数据与生成式 AI 融合起来，又避免把感知与语言坍缩成一个无监督模型。

\begin{figure}[htbp]
\centering
\includegraphics[width=\linewidth]{rag_workflow_topology.pdf}
\caption{面向立面问题的拓扑感知检索：带模态标签的原子观测聚合为资产护照、纳入层级站点索引，再经意图解析、$k$ 跳邻域扩展与上下文打包，结合带版本手册进入 LLM 生成。PDF 切块之间的嵌入邻域不能替代 BIM 邻接。}
\label{fig:ch8-rag}
\end{figure}

图~\ref{fig:ch8-rag}是多模态答案背后的检索脊柱：RGB 与红外观测在护照内仍是分卡；拓扑扩展决定哪些邻域宿主进入证据包；引用或拒答再把每个被点名的 Asset-ID 绑定到该包。

对立面运维，融合模态时不应抹掉冲突。常用查询类型包括局部查找、资产健康摘要、传播/根因阅读，以及楼层或立面尺度审计。最小装配是资产~$A$ 的护照、其邻域，以及可选的带版本手册。答案中出现的每一个 Asset-ID 或缺陷 ID 都必须落在该检索集合中，否则拒答或升级。当 RGB 与红外观测冲突时，披露优于静默平均：仅 RGB 证据不应变成地下/内部断言；仅红外证据应保留不确定标记。高分辨率立面传感与物理量化工作说明这些标记为何重要：过度宣称几何的语言会继承感知误差~\cite{zhao2024facade,zhao2026physical}。教学上实用的融合规则因此是记录级而非像素级：观测保留模态标签；由 LLM 叙述一致或冲突；在没有人工核验标记前，绝不把冲突坍缩成单一自信缺陷类型。

实现原则很直接：呈现层可以采用聊天式界面，但证据装配、持久化和孪生写入都须经由类型化服务与 harness 门禁。

当巡检工作需要多步流程时，智能体才进入：开工单、为第2--3章提议复访多边形、请求机器人补拍，或在批准后暂存孪生更新~\cite{yao2023react,schick2023toolformer,wang2024surveyagents,luo2025llmdtrobot}。记忆必须预先设计进流程：短对话状态跟踪当前问题，检索记忆刷新手册与观测卡，孪生记忆提供 Asset-ID 与拓扑，过程记忆记录已跑工具与批准~\cite{rawat2025llmmas,gao2026ifcagent,luo2025llmdtrobot}。智能体综述强调规划、记忆与工具，并警告无约束自治不安全；图孪生、人--DT、IFC 工具与运维机器人系统则把同一警告落到类型化工具与批准上~\cite{wang2024surveyagents,pan2026graphdtgpt,lu2026highwaymas,gao2026ifcagent,luo2025llmdtrobot}。因此工具增强推理应作为巡检语言系统的预期组件，而不是可有可无的插件~\cite{yao2023react,schick2023toolformer}。

对立面巡检，\emph{有界智能体}实现在第\ref{sec:ch8-llm-foundations}节的模型 harness 之内。欢迎多步检索与工具调用；最终主张仍须引用锁定到检索证据集；调度与写孪生必须经过 geofence 与人工门。小型类型化工具面，取护照、扩邻域、列观测、起草报告、提议复访，优于自由 SQL 或不受限代码执行。这一立场把通用智能体文献专门化到巡检：孪生已经编码智能体不得轻易改写的世界状态；第2--3章已经拥有飞行规划；语言智能体只作提议，不指挥飞机。近年孪生交互 MAS 研究强化同一分工：智能体编排功能、降低界面摩擦，但不取代领域批准与章节权属~\cite{lu2026highwaymas,pan2026graphdtgpt,ma2026tenquestions}。真正让这些边界在每次运行上可强制执行的是 harness。聊天记录本身不能充当门禁~\cite{yao2023react,wang2024surveyagents,liang2023helm}。

\begin{table}[htbp]
\centering
\caption{巡检语言智能体的有界工具面。}
\label{tab:ch8-tools}
\begin{tabular}{llp{5.0cm}}
\toprule
工具 & 返回 & 门禁 \\
\midrule
读取护照 & 资产摘要 + 观测 ID & 只读 \\
扩展邻域 & 拓扑列表 & 只读 \\
列出观测 & 带模态标签卡片 & 只读 \\
起草报告 & schema 约束字段 & 人工审阅 \\
提议复访 & 给第2--3章的规划提示 & 人工 + geofence \\
暂存孪生写入 & 待写入更新 & 强制批准 \\
\bottomrule
\end{tabular}
\end{table}

\begin{figure}[htbp]
\centering
\begin{tikzpicture}[node distance=7mm and 8mm]
  \node[ch8layerband, text width=10.4cm, align=center] (band)
    {L2 证据接地 $\cdot$ L3 有界流程 $\cdot$ harness 门禁};
  \node[ch8boxalt, minimum width=2.25cm, below=8mm of band.west, anchor=west] (q)
    {用户\\问题};
  \node[ch8boxrag, minimum width=2.4cm, right=of q] (a) {有界\\智能体};
  \node[ch8boxgate, minimum width=2.4cm, right=of a] (g) {人工/\\geofence 门};
  \node[ch8box, minimum width=2.4cm, below=of a] (tw) {孪生记忆\\身份锁定};
  \node[ch8boxok, minimum width=2.4cm, below=of g] (w) {批准写入/\\复访提示};
  \draw[ch8arr] (q) -- (a);
  \draw[ch8arr] (a) -- (g);
  \draw[ch8arr] (a) -- node[ch8tiny, left] {读} (tw);
  \draw[ch8arr] (tw) -- (a);
  \draw[ch8arr] (g) -- node[ch8tiny, right] {写路径} (w);
  \node[ch8defnA, below=7mm of tw.south west, anchor=north west] (d1)
    {\textbf{身份。}答案或工单中点名的宿主必须带稳定 Asset-ID。};
  \node[ch8defnB, right=2.2mm of d1] (d2)
    {\textbf{读路径。}护照、邻域与观测可自由取回。};
  \node[ch8defnC, right=2.2mm of d2] (d3)
    {\textbf{写路径。}孪生更新与复访须经批准门。};
\end{tikzpicture}
\caption{带显式身份与门禁层的有界智能体。智能体可自由读孪生记忆；不可逆动作只能经 harness 门离开，不能仅靠提示词约束。}
\label{fig:ch8-agent-gate}
\end{figure}

微型示例。经理要求在红外提示潮湿但 RGB 尚不结论处复访。有界智能体可列出带模态标签的观测，识别冲突，起草标记需人工复核的报告，并发出粗粒度复访提示。它不得在散文里虚构航迹，也不得在未批准时把新缺陷类型写入孪生。这就是本书意义上的多模态融合：模态可知的记录、接地的语言，以及尊重上游章节权属的工具。

多模态融合也与航次设计相互作用。若红外与 RGB 不同日采集，时间戳必须在观测卡片中可见；否则 LLM 可能把它们当作同时证据。若某一模态检测置信度系统性更低，信任标记应继承该不对称，禁止静默平均~\cite{zhao2024facade,zhao2026physical}。因此语言时刻的融合与感知元数据质量不可分。垃圾输入仍会得到流畅垃圾输出，只是现在带着看起来官方的引用。

智能体设计应有意保持``无聊''。在一个循环里改写 BIM、调度飞机并邮件通知承包商的壮观自治栈，不是巡检安全的好教学靶。带显式门禁的有界工具，更符合负责任运维已有的分工：分析员提议，资深人员批准，现场团队执行~\cite{luo2025llmdtrobot,gao2026ifcagent,zhang2025llmagentschema}。多智能体角色拆分为检索、拓扑扩展、报告起草与行动提议在教学上可以有用，前提是共享同一孪生记忆与同一批准政策~\cite{wang2024surveyagents,li2024surveymultiagent,rawat2025llmmas}。为同一板块发明不同 Asset-ID 的``多样智能体''不是多样性，而是失效模式。

过程记忆需要具体的巡检读法。当智能体已扩展邻域并列出观测后，不应在循环中无尽重扩。当人工拒绝某次孪生写入后，该拒绝应在会话中被记住，以免重复提交同一变更。当已提议复访时，规划提示应携带证明其必要的观测 ID，使第2--3章收到证据而非感觉~\cite{yao2023react,rawat2025llmmas}。此时记忆即运维纪律。



演示智能体时，应把工具轨迹当作可检查对象，并与答案并排展示：取了哪些护照、扩展了哪些邻域、钉扎了哪版手册、哪道门挡住了写入。轨迹使有界智能体可检查。LLM 智能体综述强调规划与工具；巡检实践应强调这些工具留下的审计痕迹~\cite{wang2024surveyagents,li2024surveymultiagent,yao2023react,schick2023toolformer}。没有轨迹的漂亮答案是演示；带轨迹的朴素答案才是运维交付物。

拆分检索、拓扑、起草与行动等多智能体角色可以澄清软件模块，但不得制造互相竞争的 Asset-ID 现实。所有角色读同一孪生记忆；只有受门禁的写入者可暂存更新；只有人类可批准。把图记忆与向量记忆结合的施工多智能体框架，是共享记忆的有用先例。对缺陷做无监督``委员会表决''不被允许~\cite{rawat2025llmmas,gao2026ifcagent,luo2025llmdtrobot}。



一次从传播问句到工单的轨迹走读足以钉牢上述纪律。把关于示意宿主 W-3-17 渗水路径的问题解析为传播意图；读取护照；扩展一跳邻域；列出带模态标签的观测；若缺上层证据则拒答；在引用或拒答规则下起草假说。若经理要求行动，则起草报告并在 RGB/IR 冲突时标记需人工复核，且仅对冲突资产提议复访并绑定观测 ID；在过程记忆中等待批准；批准后才暂存孪生写入，驳回则记录并停止。记录每一步才是产品；为了显得聪明而隐藏步骤是反模式~\cite{yao2023react,wang2024surveyagents,luo2025llmdtrobot,gao2026ifcagent}。

可把简短审计轨迹记为：意图=传播；所用工具=读取护照、扩展邻域、列出观测；检索包宿主=W-3-17、W-4-17、W-3-16；W-3-17 存在模态冲突；报告草稿待人工复核；复访待批；孪生写入=无。

工具设计还必须守住 geofence 与章节权属。复访提议必须发出规划提示，不得发出可执行无人机指令。坐标应是粗多边形或资产集合，由第2--3章在自有安全栈下细化。同样，机器人补拍请求在设施具备时也应是工单，禁止 LLM 直接打电机~\cite{luo2025llmdtrobot}。有界智能体部分是本书章节之间的组织边界。

对行动而言，状态机优于自由闲聊。 一旦用户要求工单或复访，编排器可进入显式状态机：DRAFT $\rightarrow$ REVIEW $\rightarrow$ APPROVED/REJECTED $\rightarrow$ SCHEDULED。LLM 在状态内填字段；不能因用户着急而跳过 REVIEW。强调批准的施工/运维智能体论文与这种状态机读法一致~\cite{gao2026ifcagent,zhang2025llmagentschema,rawat2025llmmas}。

不飞行也能测智能体。 智能体测试可离线：桩工具返回夹具；断言永不发生违禁写入；断言复访提示总携带观测 ID；断言被拒草稿在无新证据时不重现。这些测试比飞行航次便宜，且能更早抓住最坏失败~\cite{wang2024surveyagents,li2024surveymultiagent}。



操作员可把一页融合决策表贴在控制台旁。
\begin{center}
\begin{tabular}{ll}
\hline
RGB & IR $\rightarrow$ 语言动作 \\
\hline
有缺陷 & 有缺陷 $\rightarrow$ 双引；诊断保持审慎 \\
有缺陷 & 正常/缺失 $\rightarrow$ 披露不对称；禁止地下断言 \\
正常/缺失 & 异常 $\rightarrow$ 不确定；需人工复核 \\
时间戳冲突 & $\rightarrow$ 披露非同时性 \\
\hline
\end{tabular}
\end{center}
\vspace{0.5em}
操作员可以把这张表贴在控制台旁。它把第\ref{sec:ch8-multimodal}节变成肌肉记忆，而不需要第二套感知栈~\cite{zhao2024facade,zhao2026physical}。

何时不要调用智能体。若用户只问某资产登记了什么，读取护照与列出观测即足够，不必上工单机器。智能体只在流程问题时赚回复杂度。过早智能化会增加门禁表面积，却不增加接地~\cite{yao2023react,wang2024surveyagents}。

过程记忆的最小字段。 建议记录：session\_id、已调用工具列表、pending\_writes、rejected\_writes、approved\_revisit\_ids、last\_package\_hash。字段少才能坚持记；记了才能阻止重复提交与包被静默替换~\cite{rawat2025llmmas,gao2026ifcagent}。

人机共用的``冲突资产''队列。 所有 human\_check=true 的草稿进入同一队列，供资深巡检员按吊篮日筛选。队列不是惩罚模型，而是设施工作流。没有队列，冲突标记只是 JSON 里的装饰位~\cite{chen2023facility,luo2025llmdtrobot}。




多模态与有界智能体的教学还可以再用一条``冲突资产周''叙事收束。设西立面一周内出现三处 RGB/IR 冲突：系统不得把它们平均成单一缺陷类型，而应在吊篮名单上默认降权，并在复访提示中绑定具体观测 ID 交给第2--3章~\cite{zhao2024facade,zhao2026physical,luo2025llmdtrobot}。经理若要求``先标成已修再补证据''，harness 必须拒绝直接写入，只能暂存并等待人工与复查证据~\cite{gao2026ifcagent,lu2026highwaymas}。演示智能体时，务必并排展示 TRACE：取了哪些护照、扩了哪些邻居、钉了哪版手册、哪扇门挡住了写入。没有 TRACE 的漂亮答案仍是演示；有 TRACE 的朴素答案才是运维制品~\cite{wang2024surveyagents,liang2023helm,hellin2025bimmas}。角色拆分可以澄清软件模块，但不得制造多个 Asset-ID 真相；图孪生与人--DT MAS 的可借鉴处是共享记忆与批准。无门委员会投票决定缺陷不被允许~\cite{pan2026graphdtgpt,lu2026highwaymas,rawat2025llmmas}。



概括多模态融合与有界智能体：模态标签随观测卡片旅行，TRACE 使工具使用可检查，写入门禁保持孪生权威~\cite{zhao2026physical,yao2023react,luo2025llmdtrobot,lu2026highwaymas}。下一节把接地草稿变成设施团队可审计的报告 schema，并可选导出到既有工作系统。

\section{巡检报告生成与管理}
\label{sec:ch8-reports}

巡检答案只有落入设施管理系统已经理解的交付物时才真正有用。自由段落难以审计。实用报告 schema 可分组为身份字段，含资产与缺陷 ID、诊断、传播说明、建议行动与信任标记。身份字段必须是检索证据集的子集，这是报告层面的约束解码，与 BIM/规范助手同类~\cite{zheng2023bimassistant,joffe2025buildingcodes}。诊断字段应回指可在孪生界面打开影像的观测 ID；行动字段应可被排程系统直接阅读，而不再需要另一个 LLM 重解析散文。

\begin{table}[htbp]
\centering
\caption{巡检报告草稿的约束字段示例。}
\label{tab:ch8-schema}
\begin{tabular}{llp{6.0cm}}
\toprule
组 & 键 & 巡检经验规则 \\
\midrule
身份 & asset\_ids, defect\_ids & $\subseteq$ 检索到的孪生证据 \\
诊断 & type, severity, hypothesis & 由观测支持 \\
传播 & neighbour path & 来自拓扑扩展 \\
行动 & repair, priority & 决策支持，非自动驾驶 \\
信任 & 置信度、人工复核标记 & 模态冲突时升级 \\
\bottomrule
\end{tabular}
\end{table}

用一次走读把表~\ref{tab:ch8-schema}钉牢。设渗水路径题的检索包已含 \texttt{W-3-17} 与 \texttt{W-4-17} 护照、观测 ID \texttt{obs-2403-118} 对应 RGB 渍迹，\texttt{obs-2403-204} 对应红外异常，以及含上层宿主的一跳邻域列表。守 schema 的草稿便只点名这些 Asset-ID，把两条观测 ID 挂在诊断证据下，把传播写成向上路径，按可达性给出局部探查优先级，并因 RGB/红外不一致而抬高人工复核标记。失败草稿可以同样流畅：发明 \texttt{W-3-18}、声称``已确认渗水''，或在没有观测字段时写严重度。验收是机械的：草稿中每个身份必须是检索包的子集，每个证据 URI 必须打开到正确裁剪，冲突标记必须进入信任块~\cite{ji2023,zheng2023bimassistant,joffe2025buildingcodes}。表~\ref{tab:ch8-report-accept} 把该验收写成交接前清单。

\begin{table}[htbp]
\centering
\caption{报告草稿离开语言层之前的验收检查。}
\label{tab:ch8-report-accept}
\begin{tabular}{llp{4.6cm}}
\toprule
检查项 & 通过条件 & 坏味道 \\
\midrule
ID 子集 & 资产/缺陷 ID $\subseteq$ 证据包 & 虚构板块编码 \\
可打开证据 & 每个 URI 打开正确裁剪 & 引用表演 \\
冲突诚实 & 保留 RGB/IR 分歧 & 静默平均诊断 \\
行动可读 & 排程器可直接用字段 & 仅散文工单 \\
写入暂存 & 孪生更新仅 pending & 草稿按钮直接改孪生 \\
\bottomrule
\end{tabular}
\end{table}

\begin{figure}[htbp]
\centering
\includegraphics[width=\linewidth]{deterministic_ablation_exact_hit.pdf}
\caption{说明读者应要求的检索接地结果：在冻结孪生上，带一跳扩展的资产护照比仅切块的封装更能使答案与被引 ID 留在检索证据中。图示用于说明报告草稿应继承的证据纪律，不用于比较系统优劣。}
\label{fig:ch8-ablation}
\end{figure}

图~\ref{fig:ch8-ablation}说明读者应检查的接地结果：护照封装与受控一跳扩展能让答案 ID 和被引 ID 更贴近检索证据；跳数因此属于证据装配策略，不是自由提示旋钮~\cite{lewis2020,jeon2025hybridrag,lee2024ragsafety}。报告草稿继承同一规则：每个被点名的 Asset-ID 必须属于检索返回集合，每条证据 URI 必须可打开~\cite{ji2023,peng2024}。现场身份卫生属于本章的核心运维纪律。对于竖向栈渗水问题，系统应返回邻域护照和可打开影像的观测 ID，并在证据充足时按 schema 起草维修简报；证据不足时则拒答。需要复访时，行动提议者只为第2--3章给出规划提示。

报告管理不止一次性生成。草稿应版本化；人工签核应被记录；已批准事实可以写回第\ref{sec:ch8-llm-foundations}节引入的长期孪生记忆。被拒草稿不应静默污染护照。若机构已使用 CMMS 或 BIM 问题追踪，上表 schema 可作为可选导出映射到既有字段。也不另起一套散文 PDF 宇宙。自然语言查询界面可以帮助非专业人员提问，但权威交付物仍是约束报告加上可打开证据。实用状态机把纪律变成可强制执行的顺序：DRAFT $\rightarrow$ VALIDATE $\rightarrow$ REVIEW $\rightarrow$ APPROVED/REJECTED $\rightarrow$ 可选 EXPORT $\rightarrow$ 可选 TWIN\_WRITE。LLM 只在 DRAFT 内填字段；表~\ref{tab:ch8-report-accept} 负责 VALIDATE；人类负责 REVIEW；EXPORT 仅服务于现场已有的工作系统；只有批准路径可调用 \texttt{stage\_twin\_write}。ReAct 式工具循环与 schema 引导的 AEC 智能体已表明：多步工具使用有价值，前提是每一步仍类型化且有门~\cite{yao2023react,gao2026ifcagent,luo2025llmdtrobot}。不得因十分钟后开会就用流畅度跳过 VALIDATE~\cite{li2024surveymultiagent,rawat2025llmmas,hellin2025bimmas}。软件角色可为清晰度而拆分为检索、起草、可选映射与行动提议，但不得制造多个 Asset-ID 真相：所有角色读同一孪生记忆与同一批准政策~\cite{gao2026ifcagent,luo2025llmdtrobot,quan2026bimfacrag,buga2025llmfmsql,li2024surveymultiagent,wang2024surveyagents}。

\begin{figure}[htbp]
\centering
\includegraphics[width=\linewidth]{analysis_interface_nlq.pdf}
\caption{自然语言分析界面示例：巡检员提问后得到点名资产、可打开证据与结构化后续建议。界面布局当作交互模式阅读；验收标准仍是稳定 Asset-ID、可打开引用、拒答能力与有门禁的写入。}
\label{fig:ch8-nlq-ui}
\end{figure}

图~\ref{fig:ch8-nlq-ui}只有与第\ref{sec:ch8-dt-to-language}--\ref{sec:ch8-multimodal}节合读才有用：结构化孪生记忆、多模态证据装配与有界工具。没有这些契约的流畅聊天面板，仍只是演示。对界面模式做一次短走读，可以把区别钉牢。设巡检员问：三层以上西立面哪些宿主出现潮湿相关症状，且证据可打开。守契约的答案会点名已在检索护照集合中的 Asset-ID，列出可在孪生浏览器打开裁剪的观测 ID，在存在 RGB/IR 冲突时显式披露，并给出``扩展邻域''或``起草 schema 约束简报''一类结构化后续。失败界面在屏幕上仍可能很像：它返回流畅段落、发明板块绰号，并提供不经门禁就写孪生的``标为已修''按钮。因此读图~\ref{fig:ch8-nlq-ui}时，应沿用采购评审的同一套四项测验：稳定 Asset-ID、可打开引用、证据包偏薄时拒答、写入有门~\cite{ji2023,zhang2025,buga2025llmfmsql}。

报告 schema 应尽早与设施利益相关方协商。没人会在 CMMS 里打开的字段不会被维护；律师要求的责任字段可能需要比内部分诊更严的引用或拒答检查。教学折中是核心 schema 含身份、诊断、证据链接与信任，加上可选的组织扩展，含成本代码、承包商 ID 与可达时间窗。LLM 仅可从检索上下文或人工表单起草扩展字段；不得虚构成本代码~\cite{zheng2023bimassistant,joffe2025buildingcodes}。

人在环不是认输，而是产品。签核记录应保存谁批准了什么、批准时可见哪些证据 ID，以及是否覆盖了模态冲突标记。这些记录成为下一航次的孪生出处。没有它们，下一季的 LLM 会把上一季的推测草稿当作事实；身份卫生正是纵向语言推理的前提~\cite{zhang2025}。

同一套契约也会很快暴露薄弱演示：只有 PDF RAG 而无 Asset-ID、聊天界面打不开被引影像、未批准草稿自动写孪生、无几何归一化的严重度语言，或用嵌入邻居而非 BIM 邻接讲传播故事。命中两项及以上，系统仍只是原型。交付本身有固定一半与可选一半。固定一半是带可打开证据的 schema 约束草稿：身份、诊断、信任标记，以及打开正确裁剪的 URI。Asset-ID 卫生与引用或拒答属于这一半，但并不是产品的全部。叙述必须由已填字段渲染；不得事后从散文刮 ID；夜间链接检查应与检测器质控并列~\cite{zheng2023bimassistant,joffe2025buildingcodes,zhang2025,gao2026ifcagent}。可选一半，是把这些字段导出到现场\emph{已经在用}的工作系统。接收端可以是 CMMS、BIM 问题追踪、承包商门户，或一张简单电子表；本章并不发明新的 CMMS 产品。若存在导出，映射应短且机械：\texttt{asset\_ids} 对台账编码，证据 URI 对附件，优先级用设施量表而非 LLM 自定量表，\texttt{human\_check} 作为阻断自动关闭的挂起标记~\cite{chen2023facility,krutli2025fmrags,buga2025llmfmsql}。没有 CMMS 的现场，只要保住约束草稿、可打开证据与人工签核，仍然算成功。

交接包因此应携带 Asset-ID 词典、观测 schema、手册版本钉扎、批准日志，以及用于语言评价的冻结孪生快照。仅当已知接收系统时，才附上导出映射。报告生命周期的厚度来自这些契约：草稿先通过 ID 子集校验再进人工审，驳回进入过程记忆，批准字段可尽量免重写地导出，Asset-ID 冲突留在草稿侧栏可见~\cite{buga2025llmfmsql,chen2023facility,quan2026bimfacrag,hellin2025bimmas}。已发表的运维机器人与 IFC 智能体闭环，即语言请求、结构化动作、批准下更新孪生或模型走的是同类圆圈；立面运维把圆圈专门化，使每份可行动制品在预算或脚手架移动之前仍引用可打开的观测裁剪~\cite{luo2025llmdtrobot,gao2026ifcagent,saka2024gptconstruction}。无法打开证据的``正式报告''因此比聊天幻觉更危险~\cite{ji2023}。

每次主要飞行周结束后，冻结孪生快照，在含拒答的金标题集上重跑表~\ref{tab:ch8-report-accept}，并在更换任何托管 LLM 之前统计虚构 ID 率与断链证据率~\cite{liang2023helm,zhang2025}。跨航次主张，例如``比上季更严重''要求证据包含既往观测 ID；否则 harness 必须强制拒答，或把叙述留给仅人类填写的字段。这种季节性卫生保护纵向语言，正如 Asset-ID 卫生保护单次草稿~\cite{chen2023facility,krutli2025fmrags}。权威交付物仍是带可打开证据与人工签核的 schema 约束草稿；表~\ref{tab:ch8-report-accept} 检查这一核心，任何工作系统导出都只是可选的后续配方。下一节用实现无关准则，追问这些制品是否真的改变排程。

\section{决策支持与预测性维护}
\label{sec:ch8-dss}

决策支持在章末追问：接地语言系统究竟为了什么。运维上的答案是受控的生成：自动支撑综合报告、对历史维修记录的自然语言查询，以及能加速排程、又仍处于专业审阅下的预测性维护建议。教学上，这意味着点名资产、可打开证据、下一周吊篮的排序分诊，以及预后证据不足时的诚实拒答。评价准则因此必须可跨实现迁移。对巡检语言系统，应追问：被引 ID 是否落在证据集中~\cite{ji2023,shuster2021}；纳入既往观测与邻域后，传播题与历史记录题是否改善；面向规范/手册的 RAG 是否优于对带版本语料的纯参数回忆~\cite{lee2024ragsafety,joffe2025buildingcodes,krutli2025fmrags}；输出能否驱动运维排程；智能体工具轨迹是否遵守批准与 geofence~\cite{luo2025llmdtrobot,rawat2025llmmas}。表~\ref{tab:ch8-eval} 将这五项写成适用于不同实现的共享量规。

\begin{table}[htbp]
\centering
\caption{巡检语言系统的实现无关的评价准则。}
\label{tab:ch8-eval}
\begin{tabular}{llp{4.8cm}}
\toprule
准则 & 要问的问题 & 坏味道 \\
\midrule
接地 & 被引 ID 是否在证据中？ & 流畅但无链接散文 \\
拓扑 & 邻域是否帮助多跳？ & 仅 PDF 切块 RAG \\
语料忠实 & 是否优于参数回忆？ & 手册版本错误 \\
可运维 & FM 能否据此排程？ & 无法落地的段落 \\
智能体安全 & 写/调度是否有门？ & 静默篡改孪生 \\
\bottomrule
\end{tabular}
\end{table}

量规应按行动层计分，而不能只按句子层计分。一段从不改变维修排程的接地段落，运维价值有限；一份带可打开证据的 Asset-ID 简短排序，可能已足够安排一周吊篮。月度看板因此可跟踪虚构 ID 率、可打开证据率、模态冲突披露率、复访提示采纳率，以及在存在工作系统导出时的免重写导出率~\cite{lee2024ragsafety,luo2025llmdtrobot,rawat2025llmmas,liang2023helm}。若更换托管 LLM 提升了流畅度却恶化了虚构 ID，应拒绝替换；若流畅度不变而 TRACE 完整度上升，应接受 harness 变更。模型参数量排行榜次于这些可运维曲线。

图~\ref{fig:ch8-ablation}提供一项紧凑的教学示意：护照封装与受控一跳扩展可改善接地。应按意图与最终行动的安全性评价该结果，包括引用或拒答校验器能否拦截无证据的 ID~\cite{ji2023,peng2024,lee2024ragsafety}。

因此本章的预测性维护是三项交付，而不是一句口号。其一是\emph{分诊}：带可打开证据的 Asset-ID 排序，服务下一周吊篮。其二是\emph{条件建议}：如``若上层邻居缺少红外，则风险上升''，点名缺失证据并止步于虚构。其三是\emph{外部预后}：若存在非 LLM 时序模型，harness 可把其只读输出与被引观测并排呈现；语言系统不得从单次航次发明未来一年的剥落概率~\cite{gao2026ifcagent,zhang2025llmagentschema,zhao2026physical}。

部署视图用两张互补地图收束。图~\ref{fig:ch8-platform}给出可承载 harness 的软件分层：上层是呈现与自然语言服务，中层经 API 网关进入孪生读取与报告起草应用服务，下层是 Asset-ID、观测与带版本手册的数据存储。图~\ref{fig:ch8-transfer}再给出超越立面缺陷的巡检迁移：同一观测--几何--依赖脊柱可通过更换实体语义与边类型，实例化为建筑运维、改造排序或施工进度监测。schema 复用不是预测技能，平台图也不是第二套检测器~\cite{gao2026ifcagent,zhang2025llmagentschema,zhao2026physical}。

\begin{figure}[t]
\centering
\includegraphics[width=\linewidth]{platform_architecture.pdf}
\caption{巡检语言服务的参考平台分层：呈现与查询界面、API 网关、强制执行 harness 的应用服务，以及把 Asset-ID、观测与手册留在模型权重之外的数据存储。宜读作本章契约的部署地图，而不是强制产品蓝图。}
\label{fig:ch8-platform}
\end{figure}

\begin{figure}[t]
\centering
\includegraphics[width=\linewidth]{framework_generalizability_concept.png}
\caption{巡检迁移图：检测画像 / 实体画像 / $k$ 跳检索脊柱可从立面运维迁到改造排序与施工进度监测，关键是更换实体语义与边类型。harness 模式随脊柱迁移；领域标签不随模型权重迁移。}
\label{fig:ch8-transfer}
\end{figure}

无论读者使用手机、平板或桌面界面，同一测验都适用：被点名的 Asset-ID 必须一致，证据必须可打开；虚构板块编码或无链接楼层计数都不符合表~\ref{tab:ch8-eval} 的量规。

邻近的 2024--2026 孪生交互与 BIM 语言系统给出同一控制教训，无需第二次文献巡礼：把结构放在权重之外，用 LLM 编排，结构化路径缺失时拒答，并把治理当作设计变量~\cite{pan2026graphdtgpt,gao2026ifcagent,luo2025llmdtrobot,hellin2025bimmas,krutli2025fmrags,ma2026tenquestions}。只有表~\ref{tab:ch8-eval} 五项仍为绿灯时，用户体验收益才算数。

一条端到端走读足以把准则接到实践。设第2--7章已把为期两天的仓立面航次变成决策可用孪生：第1天完成 RGB 覆盖，第2天对疑似潮湿区补红外；感知写带模态标签观测，配准与 GeoBIM 保持宿主 Asset-ID 与设施台账对齐。第3天经理提出四类会议问题。局部查找只加载 \texttt{W-3-17} 的护照并返回可打开观测；跨航次比较在缺历史观测 ID 时拒答；传播题扩展上/侧邻域、保留 RGB/IR 标签，要么带引用起草路径假说，要么拒绝；立面审计应用第7章空间过滤，输出 schema 约束的吊篮名单。仅当经理要求行动时，有界智能体才起草工单、挂接证据 ID、在模态冲突时设置人工复核标记，并为第2--3章发出复访提示。它不给承包商发邮件，不改 BIM 几何，不清冲突旗标。已批准事实写回孪生记忆；被拒草稿进入过程记忆，并不得在同一会话中重复提交。这一条故事就是负载下的章节脊柱：孪生契约、记忆阶梯、资产中心装配、多模态披露、有门工具、报告生命周期与可运维指标。

部署应沿同一脊柱推进：先冻结 Asset-ID 与邻接边，只准入合法观测卡片；在核验拓扑前仅开放局部与健康题；随后启用一跳扩展和钉扎手册；再在写入封锁下干跑有界工具面，并在更换托管 LLM 前完成冻结快照评价与 CMMS 映射~\cite{zhang2025,zhao2026physical,peng2024,lee2024ragsafety,gao2026ifcagent,liang2023helm,ma2026tenquestions,buga2025llmfmsql}。图孪生与多智能体研究可以启发其中的角色拆分，但可信身份和邻接关系仍是起点~\cite{pan2026graphdtgpt,rawat2025llmmas,lu2026highwaymas}。

分诊语言处在诊断与排程之间。示意调度器可按 $s = 0.4\,\tilde{\ell} + 0.3\,c + 0.2\,a + 0.1\,m$ 给资产打分，其中 $\tilde{\ell}$ 为几何归一化缺陷尺度，$c$ 为检测置信度，$a$ 为可达紧迫性，$m$ 为模态冲突惩罚。系数属于设施方，不属于 LLM；语言系统应暴露每个分数背后的引用项，并让人类重排，见表~\ref{tab:ch8-triage}。表面裂缝需要几何归一化长度，潮湿线索需要多模态披露，历史修补斑块在允许``是否更严重''之前需要跨航次 ID~\cite{zhao2024facade,zhao2026physical,jeon2025hybridrag}。最小约束草稿仍是承载该分诊的机器可读载体：

\begin{verbatim}
{
  "asset_ids": ["W-3-17", "W-4-17"],
  "defect_ids": ["obs-2403-118", "obs-2403-204"],
  "diagnosis": {
    "type": "possible_moisture_path",
    "severity": "medium",
    "hypothesis": "upstairs staining with mid-storey hollow"
  },
  "propagation": ["W-4-17", "W-3-17"],
  "action": {"repair": "local_probe", "priority": 2},
  "trust": {"confidence": 0.62, "human_review": true,
            "notes": "RGB/IR disagreement on W-3-17"}
}
\end{verbatim}

草稿中每个身份都必须相对检索集做子集校验；给人看的段落由字段渲染而来，不得事后再刮 ID~\cite{zheng2023bimassistant,joffe2025buildingcodes}。

\begin{table}[htbp]
\centering
\caption{语言系统可用引用解释的示意分诊项。}
\label{tab:ch8-triage}
\begin{tabular}{llp{5.0cm}}
\toprule
项 & 来源 & 语言职责 \\
\midrule
$\tilde{\ell}$ & 孪生几何 + 观测 & 禁止虚假毫米断言 \\
$c$ & 第6章检测器 & 置信度保持可见 \\
$a$ & FM 日历 / 可达性 & 不得发明日期 \\
$m$ & RGB/IR 冲突旗标 & 披露分歧 \\
\bottomrule
\end{tabular}
\end{table}

常见失效最终都落回同一套契约。缺校验器会出现虚构 Asset-ID；用嵌入邻域替代 BIM 邻接会出现假渗水路径；红外卡片丢掉不确定标记会出现``已确认渗水''；报告起草与孪生写入共用一个无门禁按钮会出现静默写孪生；手册未钉扎会出现版本折中~\cite{ji2023,peng2024,lee2024ragsafety,luo2025llmdtrobot}。共享孪生记忆、类型化工具与门禁提供了相应的控制。采购与课堂验收可索要 Asset-ID 词典、邻接边出处、观测 schema、版本钉扎、意图--装配策略、有门工具面、CMMS 字段映射、含拒答的金标题，以及成对 TRACE 文件。图~\ref{fig:ch8-memory-ladder}--\ref{fig:ch8-agent-gate} 与表~\ref{tab:ch8-harness}、\ref{tab:ch8-triage} 是支撑该要求的承重教学对象~\cite{yao2023react,hellin2025bimmas,buga2025llmfmsql}。

望向第9章，开放问题大多是社会技术性的：老旧建筑上拓扑不完整、承包商标识冲突、安全巡检工具轨迹的智能体基准稀缺，以及``看起来已接地的答案出错时谁负责''的治理问题~\cite{saka2024gptconstruction,ma2026tenquestions}。更大模型本身并不能填平这些缺口。本章给出的建设性回应与第\ref{sec:ch8-dt}--\ref{sec:ch8-reports}节保持连续：让孪生决策可用，在资产上检索，融合模态而不抹掉冲突，约束智能体，并评价设施团队真正排得上程的东西。

概括全章：第8章是全书第四阶段——新一代决策支持。它消费第4--7章，讲授即用 LLM 技术与带巡检记忆阶梯的模型 harness，把 VLM 定位为同一引用门禁下的可选裁剪阅读器，通过资产中心检索与有界智能体接地答案，并交付历史记录查询、可审计报告与可排程的预测性维护分诊。它吸收 2024--2026 年已发表 AEC 实践，又专门化到立面运维。成功的尺度不是修辞流畅，而是点名资产、引用证据，以及仍处于专业控制之下的行动。
